\documentclass[11pt]{article}

\usepackage[margin=1in]{geometry}
\usepackage[T1]{fontenc}
\usepackage[utf8]{inputenc}
\usepackage{lmodern}
\usepackage{microtype}
\usepackage{amsmath,amssymb,amsthm}
\usepackage{mathtools}
\usepackage{stmaryrd}
\usepackage{xcolor}
\usepackage{hyperref}
\usepackage{booktabs}
\usepackage{enumitem}
\usepackage{listings}

\hypersetup{
  colorlinks=true,
  linkcolor=blue!50!black,
  citecolor=blue!50!black,
  urlcolor=blue!50!black
}

\lstdefinelanguage{Lean4}{
  morekeywords={
    import,namespace,end,open,section,noncomputable,abbrev,structure,inductive,
    def,theorem,lemma,where,match,with,if,then,else,by,fun,let,in,return,do,
    simp,calc,have,show,exact,rcases,cases,constructor,intro,apply,
    deriving,instance,class,extends,private,protected,mutual,partial,
    Type,Prop,Sort,Nat,Bool,true,false,sorry,axiom
  },
  sensitive=true,
  morecomment=[l]{--},
  morecomment=[s]{/-}{-/},
  morestring=[b]"
}

\lstset{
  language=Lean4,
  basicstyle=\ttfamily\small,
  keywordstyle=\bfseries\color{blue!50!black},
  commentstyle=\itshape\color{green!40!black},
  stringstyle=\color{red!50!black},
  showstringspaces=false,
  breaklines=true,
  frame=single,
  framerule=0.3pt,
  xleftmargin=0.5em,
  xrightmargin=0.5em,
  keepspaces=true,
  columns=fullflexible
}

% Theorem environments
\newtheorem{theorem}{Theorem}[section]
\newtheorem{lemma}[theorem]{Lemma}
\newtheorem{proposition}[theorem]{Proposition}
\newtheorem{corollary}[theorem]{Corollary}
\newtheorem{conjecture}[theorem]{Conjecture}
\theoremstyle{definition}
\newtheorem{definition}[theorem]{Definition}
\newtheorem{example}[theorem]{Example}
\theoremstyle{remark}
\newtheorem{remark}[theorem]{Remark}

% Quantale operators (stmaryrd provides \bigsqcup; define \bigsqcap if missing)
\providecommand{\bigsqcap}{\mathop{\textstyle\sqcap}\limits}

\title{Process Calculi Formalized in Lean~4:\\
\large Rho, Pi, MQ, and Modal Mu-Calculus \\
with Cross-Calculus Bridges}
\author{Zar \and Oru\v{z}i\footnote{%
``Oru\v{z}i'' denotes AI-assisted collaboration
using Claude Code (Anthropic), ChatGPT, and Codex (OpenAI).}}
\date{March 2026}

\begin{document}
\maketitle

\begin{abstract}
We present a comprehensive Lean~4 formalization of five process-algebraic
systems---the $\rho$-calculus, $\pi$-calculus, MQ-calculus,
MeTTa-calculus, and modal $\mu$-calculus---together with three
cross-calculus bridges: a $\pi$-to-$\rho$ encoding with forward
simulation and backward admin reflection, an MQ-to-MORK
interoperability bridge, and a MeTTa-to-$\rho$ shared-fragment
translation with open-map self-bisimulation witness and forward
simulation.
The formalization totals over 17{,}500 lines of Lean~4.28 across 55
files with zero \texttt{sorry} and zero custom axioms.
Shared infrastructure---typeclasses for parallel composition,
restriction, structural congruence, and reflexive-transitive
closure---eliminates approximately 40\% of cross-calculus duplication.
Quantale-valued semantics for $\mu$-calculus enables instantiation to
probabilistic, fuzzy, tropical, and PLN-evidence domains.
\end{abstract}


%% ====================================================================
\section{Introduction}
\label{sec:intro}
%% ====================================================================

Process calculi provide compositional models of concurrent, communicating
systems.  In this paper we report on the formal verification, in the
Lean~4 proof assistant~\cite{Moura2021}, of five such calculi and the
bridges that connect them.

\begin{enumerate}[leftmargin=1.5em]
  \item The \textbf{$\rho$-calculus}~\cite{MeredithRadestock2005}: a
    reflective higher-order calculus where names \emph{are} quoted
    processes, supporting structural congruence, a COMM reduction rule,
    an executable engine with soundness proof, the Spice
    calculus~\cite{MeredithPresentMoment2026} for $n$-step lookahead, and
    derived replication and restriction.
  \item The \textbf{$\pi$-calculus}~\cite{Milner1999,SangiorgiWalker2001}:
    the standard calculus of mobile processes, with a full encoding into
    the $\rho$-calculus following Lybech's name-server
    approach~\cite{Lybech2022}, including forward simulation, encoding
    morphism structure, weak bisimilarity, and backward admin reflection.
  \item The \textbf{MQ-calculus}~\cite{StayMeredithMQ2026}: a quantum
    process calculus where communication \emph{is} measurement, producing
    binary outcomes governed by Born-rule probabilities.  We formalize
    backend-parametric denotational semantics with a statevector
    instantiation and an interoperability bridge to
    MORK~\cite{MeredithStayOSLF2020}.
  \item The \textbf{MeTTa-calculus}~\cite{MeredithMeTTaCalculus2024}: a
    symmetric reflective higher-order concurrent calculus with
    backchaining, formalizing the core of the MeTTa language with COMM
    and REFL rules, an executable reduction engine, a premise-IR
    adequacy bridge, and a shared-fragment translation into the
    $\rho$-calculus with open-map self-bisimulation witness.
  \item The \textbf{modal $\mu$-calculus}~\cite{Kozen1983,BradfieldStirling2007}:
    propositional modal logic with least and greatest fixed points.  We
    provide both standard Boolean semantics and quantale-valued
    semantics, and prove that the Galois connection in
    $\rho$-calculus \emph{is} diamond-box duality.
\end{enumerate}

Three bridges tie these together: the $\pi$-to-$\rho$ encoding with
simulation theorems, an MQ-MORK binary-fold bridge, and a
MeTTa-to-$\rho$ shared-fragment translation.
All proofs are checked by the Lean~4.28 kernel with zero
\texttt{sorry}~\cite{mathlib2020}.

Table~\ref{tab:summary} gives line counts; total formalization exceeds
17{,}500 lines.

\begin{table}[t]
\centering
\small
\begin{tabular}{llrr}
\toprule
\textbf{Component} & \textbf{Location} & \textbf{Files} & \textbf{Lines} \\
\midrule
$\rho$-calculus  & \texttt{ProcessCalculi/RhoCalculus/} & 11 & 4{,}954 \\
$\pi$-calculus   & \texttt{ProcessCalculi/PiCalculus/}  & 19 & 9{,}134 \\
MQ-calculus      & \texttt{ProcessCalculi/MQCalculus/}  & 10 & 1{,}433 \\
MeTTa-calculus   & \texttt{ProcessCalculi/MeTTaCalculus/} &  9 & 1{,}174 \\
Common infra     & \texttt{ProcessCalculi/Common/}      &  4 &   246 \\
Modal $\mu$-calc & \texttt{Logic/ModalMuCalculus.lean}   &  1 &   360 \\
Quantale sem.    & \texttt{Logic/ModalQuantaleSemantics.lean} & 1 & 520 \\
\midrule
\textbf{Grand total} & & \textbf{55} & \textbf{17{,}821} \\
\bottomrule
\end{tabular}
\caption{Formalization summary: 0 sorries, 0 custom axioms throughout.}
\label{tab:summary}
\end{table}


%% ====================================================================
\section{The Rho-Calculus}
\label{sec:rho}
%% ====================================================================

The $\rho$-calculus~\cite{MeredithRadestock2005} is a reflective
higher-order process calculus in which names are quoted processes:
every name $x$ is of the form $\mathsf{@}P$ for some process~$P$.
This ``reflection'' dissolves the traditional name/process distinction
and raises expressiveness questions beyond purely behavioral
formalisms (Section~\ref{sec:mu-bridge}).
The formalization comprises 11~files totaling 4{,}954~lines under
\texttt{ProcessCalculi/RhoCalculus/}.

\subsection{Syntax}

Process syntax uses a locally nameless representation.  The core
inductive type \texttt{Pattern} (in \texttt{Types.lean}, 221~lines)
includes:

\begin{lstlisting}[caption={Rho-calculus process syntax (from \texttt{Types.lean})}]
inductive Pattern : Type where
  | pNil    : Pattern
  | pPar    : Pattern -> Pattern -> Pattern
  | pInput  : Pattern -> Pattern -> Pattern -> Pattern
  | pOutput : Pattern -> Pattern -> Pattern
  | pDrop   : Pattern -> Pattern
  | pQuote  : Pattern -> Pattern
  | pVar    : Nat -> Pattern
  deriving DecidableEq, Repr
\end{lstlisting}

\noindent
The \texttt{pQuote}/\texttt{pDrop} pair implements reflection: quoting a
process produces a name, and dropping a name recovers the process.

\subsection{Structural Congruence}

Structural congruence (\texttt{StructuralCongruence.lean}, 287~lines)
is formalized as an inductive relation including alpha-equivalence,
parallel commutativity and associativity, flattening of nested parallels,
and the property that quoting respects structural congruence:

\begin{theorem}[Quote Respects SC]
If $P \equiv_{\mathrm{SC}} Q$ then $\mathsf{@}P \equiv_{\mathrm{SC}} \mathsf{@}Q$.
\end{theorem}

\subsection{COMM Rule and Reduction}

The COMM rule (\texttt{CommRule.lean}, 298~lines) formalizes the
core communication step: when an output $\bar{x}\langle y \rangle$
meets an input $x(z).P$, the name $y$ is substituted for~$z$ in~$P$.
The full reduction relation (\texttt{Reduction.lean}, 430~lines) closes
COMM under parallel contexts and structural congruence.

\subsection{Executable Engine and Soundness}

An executable reduction engine (\texttt{Engine.lean}, 427~lines)
provides a fuel-bounded decision procedure for single-step reduction.
A separate soundness proof (\texttt{Soundness.lean}, 739~lines) shows
that every engine-computed reduction is derivable in the inductive
relation:

\begin{theorem}[Engine Soundness]
If \texttt{engineStep fuel p = some q} then \texttt{Reduces p q}.
\end{theorem}

\begin{remark}[COMM and Predecessor-Finiteness]
The COMM rule erases the channel name: the variable~$n$ binding the
communication channel appears in the LHS but not in the RHS.
Consequently, for any fixed target~$q'$, infinitely many distinct
predecessors exist (one per choice of~$n$), and the
$\rho$-calculus reduction is not predecessor-finite
(\texttt{not\_global\_hPredFinite\_langReduces\_rhoCalc}).
A companion formalization~\cite{leanOSLF2025} identifies a decidable
structural sufficient condition (\texttt{LangPredFiniteSafe}) under which
predecessor-finiteness holds; the $\rho$-calculus provably fails it
(\texttt{rhoCalc\_not\_langPredFiniteSafe}).
This is the precise structural reason that global Hennessy--Milner
equivalence requires additional quotienting for the $\rho$-calculus.
\end{remark}

\subsection{Spice Calculus: Present Moment and Lookahead}

The Spice calculus~\cite{MeredithPresentMoment2026} extends $\rho$-calculus
COMM with $n$-step lookahead (\texttt{SpiceRule.lean}, 419~lines;
\texttt{PresentMoment.lean}, 923~lines).  Agents receive the set of
states reachable in ${\leq}\, n$ steps before committing to an
interaction:
\[
  \texttt{futureStates}(p, n) = \{q \mid p \leadsto^n q\}
\]
When $n = 0$ this recovers standard $\rho$-calculus COMM.  The ``present
moment'' ($n = 1$) gives agents awareness of their immediate interaction
possibilities.  Key theorems include compositional $n$-step parallel
congruence lifting (\texttt{ReducesN.par\_head},
\texttt{ReducesN.par\_any\_pos}) and
\texttt{futureStates\_par\_head}.

\subsection{Derived Replication and Restriction}

The file \texttt{DerivedRepNu.lean} (669~lines) derives replication
$!P$ and restriction $(\nu x)\, P$ as definable constructs within
the $\rho$-calculus, proving they satisfy the expected reduction and
congruence laws.


%% ====================================================================
\section{The Pi-Calculus and the Pi-to-Rho Encoding}
\label{sec:pi}
%% ====================================================================

The $\pi$-calculus~\cite{Milner1999,SangiorgiWalker2001} is formalized
in 19~files totaling 8{,}967~lines under \texttt{ProcessCalculi/PiCalculus/}.
We follow the Lybech 2022 encoding~\cite{Lybech2022} of $\pi$-calculus
into $\rho$-calculus via a name-server approach.

\subsection{Syntax and Reduction}

The $\pi$-calculus syntax (\texttt{Syntax.lean}, 85~lines) and
reduction relation (\texttt{Reduction.lean}, 405~lines) follow
standard presentations~\cite{SangiorgiWalker2001}.

\subsection{Lybech 2022 Encoding}

The encoding (\texttt{RhoEncoding.lean}, 375~lines) maps $\pi$-calculus
names to $\rho$-calculus names via a name server that mediates
between the flat name space of $\pi$-calculus and the reflective name
space of $\rho$-calculus.

\subsection{Forward Simulation}

The key correctness result is forward simulation
(\texttt{ForwardSimulation.lean}, 1{,}387~lines): if a $\pi$-calculus
process reduces, its encoding reduces (modulo administrative steps) in
$\rho$-calculus.

\begin{theorem}[Forward Simulation]
\label{thm:forward-sim}
If $P \to_\pi P'$ then $\llbracket P \rrbracket \to_\rho^* \llbracket P' \rrbracket$
(up to administrative reductions).
\end{theorem}

\noindent
The theorem name in Lean is \texttt{encoding\_forward\_single\_step\_rf}.

\subsection{Encoding Morphism Structure}

The file \texttt{EncodingMorphism.lean} (894~lines) establishes that
the encoding preserves the algebraic structure of parallel composition,
restriction, and nil:
\[
  \llbracket P \mathbin{|} Q \rrbracket \equiv \llbracket P \rrbracket \mathbin{|} \llbracket Q \rrbracket, \qquad
  \llbracket \mathbf{0} \rrbracket \equiv \mathbf{0}, \qquad
  \llbracket (\nu a)\, P \rrbracket \equiv (\nu \hat{a})\, \llbracket P \rrbracket
\]

\subsection{Weak Bisimilarity}

Weak bisimilarity (\texttt{WeakBisim.lean}, 230~lines;
\texttt{WeakBisimDerived.lean}, 149~lines) formalizes the standard
notion from~\cite{SangiorgiWalker2001}: two processes are weakly
bisimilar when every observable reduction step of one can be matched
(up to silent $\tau$-steps) by the other, and vice versa.
The definition is Prop-valued and coinductive-style: a witness is an
explicit relation $R$ that is symmetric, closed under forward and
backward lifting, and preserves barbed observability restricted to a
finite name set~$N$ (following Lybech's $N$-restricted barbed
bisimilarity~\cite{Lybech2022}).
\texttt{WeakBisimOpenMapBridge.lean} connects this to the open-maps
approach of Joyal, Nielsen, and Winskel~\cite{JoyalNielsenWinskel1996}
by showing that weak $N$-restricted bisimilarity coincides with
path bisimilarity over a \texttt{BisimulationKit} that instantiates the
generalized open-map framework from
\texttt{CategoryTheory/GeneralizedOpenMaps.lean}.

\subsection{Backward Admin Reflection}

The largest single file in the formalization,
\texttt{BackwardAdminReflection.lean} (2{,}745~lines), proves that
administrative reductions introduced by the encoding can be reflected
backward: if $\llbracket P \rrbracket \to_\rho^* Q$ via admin steps,
then $Q$ is structurally congruent to some $\llbracket P' \rrbracket$.
The file now also exposes a fully general backward-simulation endpoint over
encoded step-sources:
\texttt{weak\_backward\_full\_simulation\_encodedSC}, together with encoded-source
specializations
\texttt{weak\_backward\_full\_simulation\_encode} and
\texttt{weak\_backward\_full\_simulation\_encode\_rfStepAt}.
These theorems flatten the reflected trace decomposition into three explicit
branches: reflexive endpoint, admin-progress branch, or reflected $\pi$-step
branch with head/tail trace data.  This completes the backward direction of
the simulation argument: every administrative reduction sequence starting
from an encoded process eventually reaches an encoded image again.


%% ====================================================================
\section{The MQ-Calculus}
\label{sec:mq}
%% ====================================================================

The MQ-calculus~\cite{StayMeredithMQ2026} unifies communication and
quantum measurement: a COMM event produces a binary measurement outcome
governed by Born-rule probabilities.  The formalization
(10~files, 1{,}433~lines) covers syntax, operational reduction,
backend-parametric denotational semantics, and a cross-calculus bridge
to MORK.
This extends the communicating quantum processes
framework of Gay and Nagarajan~\cite{GayNagarajan2005} with explicit
measurement branching in the reduction semantics.

\subsection{Syntax}

The process grammar uses de~Bruijn-style wire indexing and typed gate
specifications:

\begin{lstlisting}[caption={MQ-calculus process syntax (from \texttt{Syntax.lean})}]
inductive GateOp : Type where
  | H | X | Z
  | custom (name : String)
  deriving DecidableEq, Repr

structure GateSpec where
  op : GateOp
  target : Nat
  deriving DecidableEq, Repr

inductive Process : Type where
  | MQNil  : Process
  | MQPar  : Process -> Process -> Process
  | MQNu   : Process -> Process
  | MQGate : GateSpec -> Process -> Process
  | MQOut  : Nat -> Process
  | MQIn   : Nat -> Process -> Process -> Process
  deriving DecidableEq, Repr
\end{lstlisting}

\noindent
The \texttt{MQIn} constructor takes two continuations---one for each
measurement outcome---making measurement branching explicit in the
syntax.

\subsection{Communication as Measurement: The COMM Rule}

The COMM rule (\texttt{CommRule.lean}, 87~lines) produces a binary
measurement outcome.  A central theorem provides a constructive
both-outcomes witness:

\begin{lstlisting}[caption={COMM rule with constructive both-branches (from \texttt{CommRule.lean})}]
inductive Outcome : Type where
  | zero : Outcome
  | one  : Outcome
  deriving DecidableEq

inductive CommReduction :
    Nat -> Process -> Process -> MeasurementBranch -> Prop where
  | outcome_zero (i : Nat) (p q : Process) :
      CommReduction i p q { outcome := .zero, result := p }
  | outcome_one (i : Nat) (p q : Process) :
      CommReduction i p q { outcome := .one, result := q }

theorem comm_both_outcomes (i : Nat) (p q : Process) :
    CommReduction i p q { outcome := .zero, result := p }
    /\ CommReduction i p q { outcome := .one, result := q } :=
  (.outcome_zero i p q, .outcome_one i p q)
\end{lstlisting}

\noindent
Born-style probabilities are defined separately on quantum state
vectors; the branch relation itself remains structural and
proof-friendly.

\subsection{Backend-Parametric Denotational Semantics}

The semantics (\texttt{Backend.lean}, 380~lines;
\texttt{Denotational.lean}, 229~lines) is parameterized by a backend
interface that provides gate action, fresh allocation, branch
probabilities, and post-measurement collapse:

\begin{lstlisting}[caption={Backend interface (from \texttt{Backend.lean})}]
structure MQSemanticsBackend where
  branchProb      : Nat -> Outcome -> QState -> Real
  branchProb_nonneg : forall i out st, 0 <= branchProb i out st
  branchProb_sum_one : forall i st,
    branchProb i .zero st + branchProb i .one st = 1
  applyGate       : GateSpec -> QState -> QState
  allocFresh      : QState -> QState
  collapse        : Nat -> Outcome -> QState -> QState
\end{lstlisting}

\noindent
The current \texttt{statevectorBackend} instantiation gives executable
semantics while preserving theorem-level invariants.  A key
normalization result ensures physical consistency:

\begin{theorem}[Post-Collapse Normalization]
\texttt{collapseByOutcome\_norm\_eq\_one\_of\_raw\_pos}:
if $0 < \texttt{rawOutcomeProb}(i, \mathit{out}, \mathit{st})$ then
$\|\psi'\| = 1$ after collapse.
\end{theorem}

\subsection{MQ-MORK Bridge}

MQ measurement branching is linked to MORK binary-fold branching
(\texttt{Interoperability.lean}, 37~lines) via:

\begin{theorem}[MQ-MORK Interoperability]
\texttt{comm\_nondeterminism\_iff\_mork\_binary}: the two COMM outcomes
of a measurement correspond bijectively to the two sub-results of
a MORK binary fold step.
\end{theorem}

\noindent
This ties communication-level nondeterminism in MQ directly into the
MORK three-phase execution model~\cite{MeredithStayOSLF2020}.

\paragraph{Executable witness.}
Two MM2 scripts exercise the bridge concretely on the MORK runtime.
MORK operates on a fact database: an \texttt{exec} rule matches facts
against a pattern (after the comma) and rewrites them via a template
(\texttt{(+ ...)} adds a fact, \texttt{(- ...)} removes one).
Both scripts set up identical initial facts and an identical fold rule,
differing only in which branch is selected.

\smallskip\noindent
\textbf{Script~1} (\texttt{binary\_fold\_pick0.mm2})---selects branch~0:

\begin{lstlisting}[caption={Pick outcome 0 (from \texttt{binary\_fold\_pick0.mm2})}]
; Initial facts: a pending query with two candidate sub-results
(wait q1)
(subResult q1 0 P)
(subResult q1 1 Q)

; Choose branch 0 in this run
(pick q1 0)

; Fold rule: when (wait $q), (pick $q $i), and (subResult $q $i $r)
; all match, emit (commOutcome $q $i $r) and consume the wait token
(exec (0 fold)
  (, (wait $q)
     (pick $q $i)
     (subResult $q $i $r))
  (O (+ (commOutcome $q $i $r))
     (- (wait $q))))
\end{lstlisting}

\noindent
The pattern binds \texttt{\$q}=\texttt{q1},
\texttt{\$i}=\texttt{0}, \texttt{\$r}=\texttt{P} (since
\texttt{(pick q1 0)} selects index~0).  The template emits
\texttt{(commOutcome q1 0 P)}.

\smallskip\noindent
\textbf{Script~2} (\texttt{binary\_fold\_pick1.mm2})---selects branch~1:

\begin{lstlisting}[caption={Pick outcome 1 (from \texttt{binary\_fold\_pick1.mm2})}]
; Same initial facts
(wait q1)
(subResult q1 0 P)
(subResult q1 1 Q)

; Choose branch 1 in this run
(pick q1 1)

; Same fold rule
(exec (0 fold)
  (, (wait $q)
     (pick $q $i)
     (subResult $q $i $r))
  (O (+ (commOutcome $q $i $r))
     (- (wait $q))))
\end{lstlisting}

\noindent
Now \texttt{\$i}=\texttt{1} and \texttt{\$r}=\texttt{Q}, so
the template emits \texttt{(commOutcome q1 1 Q)}.

\noindent
The two outputs---\texttt{(commOutcome q1 0 P)} and
\texttt{(commOutcome q1 1 Q)}---are the operational witnesses for
\texttt{mork\_mq\_nondeterminism\_corresponds}: the \emph{same}
fold rule, applied to the \emph{same} query with the \emph{same}
two candidate sub-results, produces \emph{either} outcome depending
solely on the branch selection.  This is precisely the structural
binary nondeterminism that \texttt{comm\_nondeterminism\_iff\_mork\_binary}
formalizes in Lean: MQ's two \texttt{CommReduction} outcomes (zero$\to P$,
one$\to Q$) biject with MORK's two \texttt{FoldPicksSubResult} choices
(\texttt{subResult0}, \texttt{subResult1}).


%% ====================================================================
\section{The MeTTa-Calculus}
\label{sec:metta-calc}
%% ====================================================================

The MeTTa-calculus~\cite{MeredithMeTTaCalculus2024} formalizes the core
of the MeTTa language as a symmetric reflective higher-order concurrent
calculus with backchaining.  Two inputs communicate by unifying their
patterns, and patterns bind variables.  Read without the variable
convention, the communication rule lets one input's bound name capture the
other input's free name.  The formalization therefore states
\textsc{Comm} with a principal, idempotent unifier and with the two inputs
opened apart, quotes every pattern variable in the continuations, computes
the rule on a locally nameless closed form, and proves each adequacy claim
on a stated fragment.

\subsection{Syntax}

The process grammar is:
\[
  P,Q ::= \mathbf{0} \mid P \mathbin{|} Q \mid
           \mathsf{for}(t \leftarrow x)\, P \mid x\mathbin{?}P
           \mid {*}x
\]
\[
  x,y ::= @(P) \qquad
  t,u ::= (P) \mid \mathtt{true} \mid \mathtt{false} \mid n \mid s \mid \mathtt{sym}
\]
where $\mathsf{for}(t \leftarrow x)\, P$ is guarded input (listen for
pattern~$t$ on channel~$x$, then continue as~$P$), $x\mathbin{?}P$ is
reflection (inspect~$P$ and emit its future state on~$x$), $@(P)$
quotes a process into a name, and ${*}x$ drops (unquotes) a name back
to a process.

\texttt{Syntax.lean} (168~lines) encodes the grammar as abbreviations of
the shared \texttt{Pattern} AST.  Pattern variables are its free
variables:

\begin{lstlisting}[caption={MeTTa-calculus syntax (from \texttt{Syntax.lean})}]
abbrev Proc := Pattern   -- processes
abbrev Name := Pattern   -- names (quoted processes)
abbrev Term := Pattern   -- terms (data values)

def pZero    : Proc := .apply "MZero" []
def pPar     (elems : List Proc) : Proc :=
  .collection .hashBag elems none
def nQuote   (p : Proc) : Name := .apply "MQuote" [p]
def pDrop    (x : Name) : Proc := .apply "MDrop" [x]
def pFor     (t : Term) (x : Name) (k : Proc) : Proc :=
  .apply "MFor" [t, x, k]
def pReflect (x : Name) (p : Proc) : Proc :=
  .apply "MRef" [x, p]
\end{lstlisting}

\noindent
Two language definitions are provided: \texttt{mettaCalc} (the
\textsc{Comm}, \textsc{Refl} and parallel rules) and
\texttt{mettaCalcCommOnly} (the \textsc{Comm}-only fragment, which the
premise contract of Section~\ref{sec:metta-premises} uses as the lookahead
of \textsc{Refl}).

\paragraph{Binding and alpha-equivalence.}
A guard $\mathsf{for}(t \leftarrow x)\,P$ binds every variable of its
pattern~$t$, in~$t$ and in~$P$; the channel~$x$ is outside its scope.
A guard quoted inside a pattern binds its own variables.
Alpha-equivalence is decided on a locally nameless closed
form~\cite{deBruijn1972,Chargueraud2012}: \texttt{close} replaces every
bound occurrence by a code for the depth of its guard and its position in
that guard's pattern, and $P \equiv_\alpha Q$ is
$\mathtt{close}\,P = \mathtt{close}\,Q$ (\texttt{AlphaEq}).
Renaming a bound variable~$y$ of a guard to~$z$ preserves the closed form
(\texttt{alpha\_alphaEq}) when $z$ occurs neither in the pattern nor in the
continuation and $y$ is not rebound inside the continuation.  The
manuscript requires only $z \notin \mathrm{FN}(P)$; without $z \notin t$
the renaming merges two variables of one pattern and changes the guard's
reducts (\texttt{alpha\_side\_condition\_needed}).

\subsection{Structural Congruence}

The manuscript's structural congruence ($\equiv$) is the smallest
equivalence closed under:
\begin{align}
  P \mathbin{|} \mathbf{0}  &\equiv P
    \tag{\textsc{Nil}} \\
  P \mathbin{|} Q            &\equiv Q \mathbin{|} P
    \tag{\textsc{Comm}} \\
  (P \mathbin{|} Q) \mathbin{|} R
                              &\equiv P \mathbin{|} (Q \mathbin{|} R)
    \tag{\textsc{Assoc}} \\
  @({*}x)                    &\equiv x
    \tag{\textsc{QuoteDrop}}
\end{align}

\texttt{StructuralCongruence.lean} (67~lines) defines the inductive
relation \texttt{SC} (notation~$\approx_m$) with 11~constructors: the
equivalence closure, a constructor \texttt{alpha} for syntactic identity,
the singleton, unit, commutativity and associativity laws of the bag,
permutation, and pointwise congruence.  \textsc{QuoteDrop} is an equation
of the language definition \texttt{mettaCalc}.  Alpha-equivalence is not a
constructor of \texttt{SC}; it is \texttt{AlphaEq}, and the corrected
relational rules below use it as their equivalence.

\begin{lstlisting}[caption={Structural congruence (from \texttt{StructuralCongruence.lean})}]
inductive SC : Proc -> Proc -> Prop where
  | refl   (p : Proc) : SC p p
  | symm   (p q : Proc) : SC p q -> SC q p
  | trans  (p q r : Proc) : SC p q -> SC q r -> SC p r
  | alpha  (p q : Proc) : p = q -> SC p q
  | par_singleton (p : Proc) : SC (pPar [p]) p
  | par_nil_left  (p : Proc) : SC (pPar [pZero, p]) p
  | par_nil_right (p : Proc) : SC (pPar [p, pZero]) p
  | par_comm  (p q : Proc) :
      SC (pPar [p, q]) (pPar [q, p])
  | par_assoc (p q r : Proc) :
      SC (pPar [pPar [p,q], r]) (pPar [p, pPar [q,r]])
  | par_perm  (xs ys : List Proc) :
      xs.Perm ys -> SC (pPar xs) (pPar ys)
  | par_cong  (xs ys : List Proc) :
      xs.length = ys.length ->
      (forall i hx hy, SC (xs.get {val:=i,isLt:=hx})
                           (ys.get {val:=i,isLt:=hy})) ->
      SC (pPar xs) (pPar ys)
\end{lstlisting}

\noindent
Paper-mapped equivalence theorems:
\begin{itemize}
\item \texttt{paper\_equiv\_par\_nil\_left}: $\mathbf{0} \mathbin{|} P \approx_m P$
\item \texttt{paper\_equiv\_par\_nil\_right}: $P \mathbin{|} \mathbf{0} \approx_m P$
\item \texttt{paper\_equiv\_par\_comm}: $P \mathbin{|} Q \approx_m Q \mathbin{|} P$
\item \texttt{paper\_equiv\_par\_assoc}: $(P \mathbin{|} Q) \mathbin{|} R \approx_m P \mathbin{|} (Q \mathbin{|} R)$
\end{itemize}

\subsection{Reduction Semantics}
\label{sec:metta-reduction}

\paragraph{The capture-avoiding \textsc{Comm}.}
Two guarded inputs on the same channel communicate when their patterns
unify:
\begin{equation}
  \frac{\sigma \text{ a principal idempotent unifier of } t \text{ and } u
        \qquad \text{the two inputs apart}}
       {\mathsf{for}(t \leftarrow x)\, P \mathbin{|}
        \mathsf{for}(u \leftarrow x)\, Q
        \;\longrightarrow\;
        P\dot\sigma \mathbin{|} Q\dot\sigma}
  \tag{\textsc{Comm}}
\end{equation}
Here $\dot\sigma$ is \emph{dot-substitution}: each pattern variable~$v$
becomes the name $@(v\sigma)$, the quotation of its value, so a variable
that~$\sigma$ leaves unbound becomes $@v$.  This \emph{uniform dot} is the
manuscript's dot of a principal binding list that names every pattern
variable, a bound one with its value and an unbound one with itself
(\texttt{completeBindings}, \texttt{applyDot\_completeBindings}).  Because
$\sigma$ is idempotent, the dot is applied once to each continuation.  The
two inputs are \emph{apart} under the variable
convention~\cite{Barendregt1984,UrbanBerghoferNorrish2007}
(\texttt{Apart}): the two patterns share no variable, neither input
mentions the other's pattern variables, and no guard inside a
continuation rebinds its own input's pattern variables.  As the convention
prescribes, bound names are renamed apart, up to alpha-equivalence, before
the rule applies.

\paragraph{The relational rules.}
\texttt{Corrected.Step admit} has the manuscript's four rules:
\textsc{Comm} with a principal unifier and an admission on the redex,
\textsc{Par}, \textsc{Equiv} by~$\equiv_\alpha$, and \textsc{Refl}, whose
premise is the same relation, so that reflection is recursive:
\begin{equation}
  \frac{P \longrightarrow P'}
       {x\mathbin{?}P \;\longrightarrow\;
        \mathsf{for}((P') \leftarrow x)\, \mathbf{0}}
  \tag{\textsc{Refl}}
\end{equation}
The admission \texttt{conventional} is \texttt{Apart}.  The admission
\texttt{uniform}, the semantics, adds that the binding list names every
pattern variable (\texttt{Complete}), so that the dot is uniform;
\texttt{uniformCalc} restricts it to guard-free patterns and to patterns and
continuations of the calculus fragment.  The admission \texttt{closing} adds
to \texttt{conventional} that both patterns are guard-free, that patterns and
continuations belong to the calculus fragment, and that the unifier binds
every pattern variable a continuation reads to a ground term.
\texttt{Principal t u bs} says that one pass of~\texttt{bs} makes $t$
and~$u$ equal and that every unifier of $t$ and~$u$ absorbs it.

\paragraph{Unification.}
\texttt{Corrected.unify} encodes patterns injectively as first-order terms
and runs a Martelli--Montanari solver~\cite{MartelliMontanari1982} with an
occurs check.  The solver is evaluated by the kernel and proved equal to a
total unification algorithm, and \texttt{unify} returns each moved
variable with its final value.  It is sound, most general and idempotent,
and it fails exactly when no substitution unifies the two
patterns~\cite{Robinson1965}.

\paragraph{The executable step.}
\texttt{cstep}~$P$ works on $\mathtt{close}\,P$.  At a binary composition
of two guards on one channel with guard-free patterns, it opens the two
inputs with distinct names fresh for the whole process, unifies the opened
patterns, completes the unifier with the opened variables it leaves
unbound, and applies the dot once to each opened continuation.  It also
steps the left component of a composition, and it steps a reflected
process with the same function, so it has the nested reflection result.

\begin{lstlisting}[caption={The corrected COMM (from \texttt{CorrectedReduction.lean})}]
def commClosed (avoid : List String) (first second : Pattern) :
    List Pattern :=
  match asGuard first, asGuard second with
  | some (t, x, k), some (u, y, l) =>
      if x = y /\ guardFree t = true /\ guardFree u = true then
        let leftNames  := freshNames avoid 0 (codeCount t)
        let rightNames := freshNames avoid (codeCount t) (codeCount u)
        match unify (openPattern leftNames t)
                    (openPattern rightNames u) with
        | some bs =>
            let all := completeBindings
              (pvars (openPattern leftNames t) ++
               pvars (openPattern rightNames u)) bs
            [par (applyDot all (openWith leftNames k))
                 (applyDot all (openWith rightNames l))]
        | none => []
      else []
  | _, _ => []

def cstep (p : Pattern) : List Pattern :=
  stepClosed (pvars (close p)) (close p)
\end{lstlisting}

\begin{theorem}[The corrected step]
\label{thm:metta-corrected}
Each claim holds on the fragment named.
\begin{enumerate}[nosep]
\item \emph{Unifier, all patterns:} \texttt{unify} returns a principal,
  idempotent unifier when one exists and fails otherwise
  (\texttt{unify\_sound}, \texttt{unify\_most\_general},
  \texttt{unify\_idempotent}, \texttt{unify\_eq\_none\_iff}); one pass of
  its dotted bindings is $P\dot\sigma$ (\texttt{applyDot\_unify}).
\item \emph{Alpha-invariance, all processes:} if $P \equiv_\alpha P'$ then
  $\mathtt{cstep}\,P = \mathtt{cstep}\,P'$
  (\texttt{cstep\_alpha\_invariant}, \texttt{cstep\_alpha\_rule\_invariant}).
\item \emph{Alpha, tame processes:} on processes of the calculus fragment
  whose guards have guard-free patterns, the manuscript's renaming rule,
  closed under equivalence and contexts, is exactly equality of closed forms
  (\texttt{alpha\_iff\_alphaEq}).  Outside them a code inside a process is
  indistinguishable from a bound occurrence
  (\texttt{alphaEq\_not\_alpha\_with\_code}).
\item \emph{Soundness, the calculus fragment:} every $Q \in \mathtt{cstep}\,P$
  is a step $P \longrightarrow Q$ of the relational rules with the uniform
  dot (\texttt{cstep\_sound\_uniform}).
\item \emph{Completeness up to alpha, the closing fragment:} for every step
  $P \longrightarrow Q$ whose \textsc{Comm} instances are closing there is
  $Q' \in \mathtt{cstep}\,P$ with $Q' \equiv_\alpha Q$
  (\texttt{cstep\_complete}, \texttt{cstep\_adequate}).
\item \emph{Completeness up to renaming, guard-free patterns:} for every
  step $P \longrightarrow Q$ of the uniform dot on guard-free patterns there
  are $Q' \in \mathtt{cstep}\,P$ and a renaming~$\rho$ that fixes every free
  name of~$P$, is injective on the other names of~$Q'$, and makes
  $\rho\,Q' \equiv_\alpha Q$ (\texttt{cstep\_complete\_residual}).  The
  renaming comes from the fact that two principal unifiers agree up to a
  renaming of their residual variables (\texttt{principal\_renaming}).
\item \emph{Parallel composition:} modulo the commutativity, associativity
  and unit laws, applied to the composition being executed and never inside
  patterns, the step that fires \textsc{Comm} at any pair of guards and
  \textsc{Refl} at any reflected component is sound for calculus processes
  (\texttt{cstepBag\_sound}), and on the closing fragment it is complete for
  the rules without \textsc{Refl} (\texttt{cstepBag\_complete}).  When
  \textsc{Refl} quotes the canonical representative of its reduct, the
  listener depends only on the class of the reduct
  (\texttt{reflect\_respects\_bagEq}), and the step is complete on the
  closing fragment for all four rules (\texttt{cstepCanon\_complete}).
\item \emph{Firing:} for any principal unifier of inputs that are apart,
  with first-order continuations and a store whose variables are apart
  from the patterns, \textsc{Comm} is one valid firing of a multiset
  rewriting step with an equation store; every solution of the resulting
  store reads each output as the uniform dot below, which is the dotted
  continuation wherever the unifier binds the variables it reads
  (\texttt{comm\_adequacy\_corrected}).
\end{enumerate}
\end{theorem}

\paragraph{Why the uniform dot.}
With the manuscript's dot of an arbitrary principal binding list, the
reduct depends on how a most general unifier is oriented: for
$\{X \mapsto Y\}$ and $\{Y \mapsto X\}$ the continuation~$X$ reads $@Y$
and~$X$ (\texttt{manuscript\_dot\_depends\_on\_orientation}), and the
reduct $X \mathbin{|} @X$ of
$\mathsf{for}(X \leftarrow @\mathbf{0})\,X \mathbin{|}
 \mathsf{for}(Y \leftarrow @\mathbf{0})\,Y$ is related by no renaming to
any executable reduct, since one of its continuations is not quoted
(\texttt{manuscript\_dot\_not\_up\_to\_renaming}).  The uniform dot reads
$@Z \mathbin{|} @Z$ for one residual name~$Z$, whatever the orientation
(\texttt{uniform\_dot\_up\_to\_renaming}).  It agrees with the manuscript's
dot wherever every read variable is bound (\texttt{uniformDot\_eq\_applyDot}),
and every solution of the firing's store reads it, for any principal
unifier (\texttt{firing\_reads\_uniform\_dot}).  The renaming of
completeness cannot be injective on all names: the relational rule may
choose as a residual name a free name of the surrounding process, which
the executable, fresh for the whole process, never does
(\texttt{residual\_meets\_context}).

\paragraph{Reflection modulo the monoid laws.}
A reflected listener quotes a process, and a partner can compare quoted
processes.  If the listener quoted whichever representative of the reduct
a derivation produced, one class of reducts would give listeners that a
partner tells apart (\texttt{representative\_reflection\_two\_listeners},
\texttt{partner\_tells\_listeners\_apart}), and the order of a parallel
composition would become observable.  In the reflective
calculi a name is a quoted process taken up to structural congruence, so
that $P \equiv Q$ implies $@P \equiv_N @Q$~\cite{MeredithRadestock2005};
order is not observable while multiplicity is.  Following that, the
formalization's \textsc{Refl} quotes the canonical representative of the
reduct's class: its closed components, sorted by an injective structural
code, with their multiplicities (\texttt{canonicalBag},
\texttt{bagEq\_iff\_canonicalBag\_eq}).  From either order of the control's
two inputs, every listener is the same one
(\texttt{canonical\_reflection\_one\_listener}).  The canonical order is
exact for ground names.  A pattern that quotes a bag with variables is
still compared by unification as written, and a partner quoting another
order misses the listener (\texttt{written\_order\_partner\_misses});
comparing such patterns modulo the laws needs unification modulo
associativity and commutativity, which yields complete sets of unifiers
rather than one, and is not part of the rules.

\paragraph{Patterns that quote guards.}
The executable step does not fire on patterns that quote guards, even when
the two patterns are alpha-equivalent
(\texttt{guard\_quoting\_partners\_refused}).  Unifying such patterns up to
their bound names is the problem of nominal unification~\cite{Pitts2003}.
On closed forms it becomes first-order: closing commutes with a
substitution of closed values that captures nothing
(\texttt{closeIn\_psubst\_closed}), so such a substitution unifies two
patterns up to alpha exactly when it unifies their closed forms, on which a
bound name is a code compared as a constant
(\texttt{alphaEq\_psubst\_iff}).  The first-order unifier of the closed
forms is principal among capture-free alpha-unifiers
(\texttt{alpha\_unifier\_absorbs}); when it binds a pattern variable to a
code that a quoted guard binds, no capture-free alpha-unifier exists
(\texttt{no\_alpha\_unifier\_of\_escaping}).  Extending the executable
\textsc{Comm} and its admission to such patterns is the remaining step.

\paragraph{The naive reading.}
\texttt{Reduction.lean} keeps a naive executable reading of the rules,
which unifies the patterns as written with a triangular binding list read
in one pass and substitutes into both continuations without renaming the
inputs apart.  Kernel-computed examples separate it from the rules:
\begin{itemize}[nosep]
\item \emph{Capture.}
  $\mathsf{for}(a \leftarrow @\mathbf{0})\,a \mathbin{|}
   \mathsf{for}(\mathbf{0} \leftarrow @\mathbf{0})\,a$ reduces to
  $@\mathbf{0} \mathbin{|} @\mathbf{0}$, identifying the second input's
  free~$a$ with the first input's bound~$a$; its alpha-variant reduces to
  $@\mathbf{0} \mathbin{|} a$, which is the corrected reduct of both.
\item \emph{A triangular list read once.}
  $\mathsf{for}(f(X,X) \leftarrow @\mathbf{0})\,X \mathbin{|}
   \mathsf{for}(f(Y,a) \leftarrow @\mathbf{0})\,\mathbf{0}$ reduces to
  $@Y \mathbin{|} \mathbf{0}$, while every unifier reads $@a$.
\item \emph{No renaming.}
  $\mathsf{for}(X \leftarrow @\mathbf{0})\,X \mathbin{|}
   \mathsf{for}(f(X) \leftarrow @\mathbf{0})\,\mathbf{0}$ has no naive
  reduct; its corrected reduct is $@f(z) \mathbin{|} \mathbf{0}$ for a
  fresh~$z$.
\item \emph{The occurs check defeated.} $f(X, g(Z), Z, X)$ and
  $f(f(Y), Y, W, W)$ have no unifier, yet the naive unifier accepts them and
  the naive step fires.
\item \emph{Nested reflection.} The naive \textsc{Comm}-only lookahead has
  no result for a reflection of a reflection, so it misses a step of the
  relational rules (\texttt{naive\_step\_misses\_nested\_reflection}).
\end{itemize}
The defect is the missing convention, not only the unifier: the rules read
literally, with no admission on \textsc{Comm}, relate the canary to the
capturing reduct even with a principal unifier and alpha-equivalence
(\texttt{literal\_rule\_admits\_capture}), and the corrected rules exclude
it under the convention alone (\texttt{capture\_not\_conventional\_step}).

\subsection{The MeTTa-Calculus as a GSLT}

\texttt{GSLTInteraction.lean} (265~lines) gives the behavioral theory
\texttt{mettaCalcGSLT}: its terms are the processes of the calculus
fragment, its equations are alpha-equivalence, and its rewrites are the
relational rules under the convention with the uniform dot.
\textsc{Equiv} makes the rewrites
respect the equations on both sides.  Every reduct of the corrected step
that is a calculus process is a rewrite
(\texttt{executable\_step\_to\_mettaCalcGSLT}); on the closing fragment
every rewrite is a reduct up to alpha (\texttt{closing\_step\_executes}),
and on guard-free patterns every rewrite is a reduct up to renaming its
residual names (\texttt{uniform\_step\_executes\_up\_to\_renaming}).
A second theory, \texttt{mettaCalcBagGSLT}, takes the monoid laws as its
equations and the rules with canonical reflection as its rewrites, so that
the order of a parallel composition is not observable; on the closing
fragment every rewrite is a reduct of the step with canonical reflection,
up to the laws (\texttt{closing\_bag\_step\_executes}).
The naive reading's capturing reduct of the canary is not a rewrite
(\texttt{mettaCalcGSLT\_no\_capture}).
A reduct may keep a guard of a continuation in closed form, and is then
not itself a term of the theory.  The authored language definition
supplies the interactive presentation: the process sort, the parallel
constructor and the symmetric \textsc{Comm} rewrite, a base interaction
whose one premise is a relation query.  The theory the generic engine
presents from the authored definition has closed ground terms as its
carrier, which contain no pattern variable, so it can express neither
binding nor capture and is not used as the behavioral theory.  The direct
contact of two guarded parties through the located-transaction layer
takes its premise as a parameter; with the corrected premise, inputs apart
and the principal unifier of their patterns completed to name every pattern
variable, every contact is a root \textsc{Comm} of the relational rules with
the uniform dot (\texttt{correctedContact\_step}).

\subsection{Premise Intermediate Representation and Adequacy}
\label{sec:metta-premises}

\texttt{Premises.lean} (95~lines) declares a \texttt{PremiseProgram}---a
datalog-style intermediate representation (IR) with two relations and
two builtins:
\[
  \mathit{mettaComm}(t, u, P, Q, P', Q') \qquad
  \mathit{mettaStepNoReflect}(\mathit{src}, \mathit{tgt})
\]
\[
  \mathit{mettaCommWitness}(t, u, P, Q) \rightsquigarrow \mathit{MRef}(P', Q')
  \qquad
  \mathit{mettaCommOnlyStep}(\mathit{src}) \rightsquigarrow \mathit{tgt}
\]
with machine-checked stratification and wellformedness proofs.

\texttt{Adequacy.lean} (757~lines) answers the contract with the corrected
\textsc{Comm}.  A rule sees only its own bag, so the names a residual must
avoid cannot come from the rule; the runtime supplies the free names of the
process being stepped, which is how the corrected step chooses its fresh
names too.

\begin{theorem}[Premise--Runtime Adequacy]
\label{thm:metta-adequacy}
The following are proven:
\begin{enumerate}[nosep]
\item \texttt{runtime\_covers\_declared\_relations},
  \texttt{runtime\_covers\_declared\_builtins}: every declared relation and
  builtin is implemented by the runtime dispatcher.
\item \texttt{comm\_witness\_adapter\_spec}: the witness for
  $\mathit{mettaCommWitness}(t, u, P, Q)$ is the corrected \textsc{Comm} of
  the two inputs, closed, opened apart with names fresh for its four
  arguments and the supplied names, unified, and dotted once.
\item \texttt{comm\_witness\_sound}: at a root redex of calculus terms, on
  any channel, every witness is a step of the relational rules under the
  convention.
\item \texttt{stepViaPremiseContract\_sound}: every reduct of the run of the
  authored rules on a calculus process is a step of the relational rules
  modulo the monoid laws.  The proof goes through the engine's relational
  specification of its fuel-bounded run and inverts the authored rules'
  matches (\texttt{engine\_sound}).
\item \texttt{stepViaPremiseContract\_complete\_comm}: if two elements of
  one bag communicate on the closing fragment, the run produces the reduct,
  the other elements kept, up to the monoid laws.
\item \texttt{premise\_contract\_comm\_canary},
  \texttt{premise\_contract\_no\_capture},
  \texttt{premise\_contract\_comm\_bagEq}: the authored rules reduce the
  canary to its corrected reduct in both orders of the bag and never to the
  capturing one, and each of their reducts is the corrected step's reduct
  up to the monoid laws.
\end{enumerate}
\end{theorem}

\noindent
Outside the fragment of the completeness item the run stops short: two
inputs in different bags have no reduct, although the monoid laws bring
them together (\texttt{nested\_bag\_missed}), and a reflection whose process
is itself a reflection has none, because the contract's \textsc{Refl}
premise is the declared \textsc{Comm}-only lookahead
(\texttt{nested\_reflection\_missed}).  Without the supplied names a
witness is fresh for its redex only, and its residual name can be a free
name of the surrounding bag (\texttt{supply\_keeps\_residuals\_fresh}).

\subsection{Interoperability with \texorpdfstring{$\rho$}{rho}-Calculus}

\texttt{Interoperability.lean} (319~lines) provides a shared-fragment
translation from the quote/drop/parallel subset of MeTTa-calculus into
$\rho$-calculus.  The translation \texttt{toRhoSharedProc?}~succeeds on
the translatable fragment and fails gracefully outside it.

For any successfully translated process, the file proves path-bisimulation
self-witnesses, intrinsic rho-side inversion, and shared-core
forward/backward simulation:

\begin{itemize}
\item \texttt{shared\_translation\_pathBisim\_self}: the translated image
  is path-bisimilar to itself under the $\rho$-calculus
  \texttt{BisimulationKit}.
\item \texttt{shared\_translation\_esBisim\_self}: the corresponding
  $(E,S)$-bisimilarity corollary.
\item \texttt{rhoSharedCore\_step\_inversion}: intrinsic inversion of
  \texttt{RhoSharedCoreReduces}.
\item \texttt{rhoSharedCore\_stepStar\_inversion}: intrinsic inversion of
  \texttt{RhoSharedCoreReducesStar}.
\item \texttt{sharedCore\_step\_forward\_restricted}: one-step forward
  simulation on the quote/drop/parallel shared-core fragment.
\item \texttt{sharedCore\_step\_backward\_restricted}: source-anchored
  one-step backward simulation.
\item \texttt{sharedCore\_step\_bisimulation}: combined one-step
  forward+backward shared-core bisimulation correspondence.
\item \texttt{sharedCore\_stepStar\_forward\_restricted}: reflexive-transitive
  forward simulation.
\item \texttt{sharedCore\_stepStar\_backward\_restricted}: reflexive-transitive
  backward simulation.
\item \texttt{sharedCore\_stepStar\_bisimulation}: combined star-level
  forward+backward shared-core bisimulation correspondence.
\end{itemize}

\noindent
Thus, forward and backward simulation are now formalized for translated
shared-core quote/drop/parallel steps (and their star closure).  Remaining
future work is to extend this from shared-core to the full MeTTa fragment
(including \textsc{Comm}/\textsc{Refl}) and prove an assumption-free
cross-language bisimulation theorem.

\subsection{Regression and Paper Map}

\texttt{PaperMap.lean} (190~lines) provides a theorem-level index mapping
each clause of the specification~\cite{MeredithMeTTaCalculus2024} to its
Lean proof.  Its reduction canaries are kernel computations of the
corrected step: the canary reduces to exactly
$@\mathbf{0} \mathbin{|} a$ (\texttt{paper\_comm\_positive}), its
reflection emits the listener that quotes this reduct
(\texttt{paper\_refl\_positive}), and mismatched patterns block both
(\texttt{paper\_comm\_negative}, \texttt{paper\_refl\_negative});
\texttt{paper\_comm\_relational} derives the same reduct by the relational
rules, and \texttt{paper\_comm\_sound}, \texttt{paper\_comm\_complete},
\texttt{paper\_comm\_alpha\_invariant} and
\texttt{paper\_comm\_unifier\_principal} index
Theorem~\ref{thm:metta-corrected}.  The index also wraps step and star
forward simulation, backward simulation, intrinsic backward inversion,
and combined bisimulation for the shared fragment
(\texttt{paper\_shared\_to\_rho\_*} theorems).
\texttt{CoreMain.lean} now also exposes an explicit HM endpoint map
(\texttt{coreMain\_hm\_endpoint\_recommendation\_map}) distinguishing global
canonical wrappers (assumption-scoped predecessor-finiteness) from the
assumption-free scoped SC-quotiented family.
\texttt{Regression.lean} (79~lines) fixes the exact outputs of the
corrected step on the positive and negative \textsc{Comm} and \textsc{Refl}
canaries and in a parallel context: as a list, as a set, and modulo the
monoid laws, where \textsc{Comm} fires at either order of the pair.
\texttt{CorrectedControls.lean} holds the separating examples against the
naive reading.
\texttt{RelationNames.lean} (28~lines) is a single source of truth for
relation and builtin string identifiers.


%% ====================================================================
\section{Modal Mu-Calculus and the Rho-Mu Bridge}
\label{sec:mu-bridge}
%% ====================================================================

The modal $\mu$-calculus~\cite{Kozen1983,BradfieldStirling2007} extends
Hennessy-Milner logic with least and greatest fixed-point operators.
Our formalization covers standard semantics
(\texttt{ModalMuCalculus.lean}, 360~lines).

\subsection{Syntax}

Formulas use de~Bruijn indices for bound variables:
\[
  \phi ::= \top \mid \bot \mid \neg \phi \mid \phi \land \psi \mid
           \phi \lor \psi \mid \langle a \rangle \phi \mid [a] \phi
           \mid \mu X.\phi \mid \nu X.\phi \mid X
\]
where $\langle a \rangle$ is diamond (possibility), $[a]$ is box
(necessity), $\mu X$ is least fixed point, and $\nu X$ is greatest
fixed point.

\subsection{Rho-Calculus as LTS}

The $\rho$-calculus reduction relation $p \leadsto q$ is a labelled
transition system (LTS) with
states = patterns, actions = $()$ (unlabeled), and transitions =
$\leadsto$.  This directly induces modal operators:

\begin{definition}[Reduction-Based Modalities]
\begin{align*}
\texttt{possiblyProp}\, \phi\, p &:= \exists q.\, p \leadsto q \land \phi\, q \quad \text{(diamond)} \\
\texttt{relyProp}\, \phi\, p &:= \forall q.\, q \leadsto p \to \phi\, q \quad \text{(box)}
\end{align*}
\end{definition}

\begin{theorem}[Galois Connection IS Modal Duality]
\label{thm:galois}
\[
(\forall p.\, \texttt{possiblyProp}\, \phi\, p \to \psi\, p) \iff
(\forall p.\, \phi\, p \to \texttt{relyProp}\, \psi\, p)
\]
This is exactly $\langle\rangle \phi \subseteq \psi \iff
\phi \subseteq [\,] \psi$.
\end{theorem}

\noindent
The proof is by direct unfolding of definitions
(\texttt{galois\_connection} in \texttt{Reduction.lean}).
This validates the OSLF construction~\cite{MeredithStayOSLF2020,leanOSLF2025}:
operational semantics naturally gives rise to modal logic.

\subsection{Spice Calculus as Iterated Diamond}

The Spice lookahead corresponds precisely to iterated modal operators:
\begin{align*}
  \texttt{presentMoment}(p) &\equiv \langle\rangle \phi \quad \text{(1-step)} \\
  \texttt{futureStates}(p, n) &\equiv \langle\rangle^n \phi \quad \text{($n$-step)} \\
  \texttt{reachableViaStarClosure}(p) &\equiv \nu X.\, \phi \lor \langle\rangle X \quad \text{(eventually)}
\end{align*}
All three equivalences hold by definitional equality (\texttt{simp}).
Precognitive agents computing Spice lookahead are computing
$\mu$-calculus temporal properties.

\subsection{Conjecture: Rho Strictly More Expressive Than Mu}
\label{sec:expressiveness}

\begin{conjecture}[$\rho$-Calculus Strictly More Expressive]
\label{conj:rho-mu-expressiveness}
There exists a property expressible in $\rho$-calculus that is not
expressible in modal $\mu$-calculus.  Specifically: $\mu$-calculus
cannot express ``this process has the capability to quote itself.''
\end{conjecture}

The intuition is that $\mu$-calculus formulas are determined entirely by
LTS behavior (which transitions occur), while $\rho$-calculus processes
can inspect their own syntactic structure via quote/unquote.  Two
processes $x[x]$ (self-quoting) and $y[z]$ (not self-quoting) are
behaviorally indistinguishable---both are dead---yet differ structurally.
No $\mu$-formula should be able to distinguish them.

A proper formalization would require: (1)~a coinductive bisimulation
relation (not just deadlock equivalence), (2)~a simulation theorem
showing $\rho$ can faithfully encode $\mu$-calculus evaluation, and
(3)~the separation argument over the actual $\rho$-calculus types from
\texttt{ProcessCalculi/RhoCalculus/}.  This remains future work.

\begin{remark}[G\"odel Analogy]
This is analogous to G\"odel's incompleteness theorem: first-order logic
cannot express ``this formula is provable''; adding a provability
predicate (reflection) gives strictly more expressive power.  Similarly,
$\mu$-calculus cannot express ``this process quotes itself'';
$\rho$-calculus can, via the quote operator $\mathsf{@}$.
\end{remark}


%% ====================================================================
\section{Quantale-Valued Semantics}
\label{sec:quantale}
%% ====================================================================

Standard $\mu$-calculus has Boolean semantics.  We generalize to
\emph{quantale-valued} semantics~\cite{Kozen1983} where satisfaction
takes values in a commutative quantale~$Q$
(\texttt{ModalQuantaleSemantics.lean}, 520~lines).

\begin{definition}[QLTS]
A \emph{quantale-labeled transition system} is a tuple
$(S, A, \mathrm{trans} : S \times A \times S \to Q)$ where $Q$ is a
commutative quantale with residuation.
\end{definition}

\begin{definition}[Quantale-Valued Satisfaction]
\begin{align*}
\texttt{qSatisfies}(s, \phi \land \psi) &= \texttt{qSatisfies}(s, \phi) \sqcap \texttt{qSatisfies}(s, \psi) \\
\texttt{qSatisfies}(s, \phi \lor \psi) &= \texttt{qSatisfies}(s, \phi) \sqcup \texttt{qSatisfies}(s, \psi) \\
\texttt{qSatisfies}(s, \langle a \rangle \phi) &= \bigsqcup_{s'} \mathrm{trans}(s, a, s') \otimes \texttt{qSatisfies}(s', \phi) \\
\texttt{qSatisfies}(s, [a] \phi) &= \bigsqcap_{s'} \mathrm{trans}(s, a, s') \Rightarrow \texttt{qSatisfies}(s', \phi) \\
\texttt{qSatisfies}(s, \mu X.\phi) &= \bigsqcap \{P \mid \texttt{transformer}(P) \sqsubseteq P\} \\
\texttt{qSatisfies}(s, \nu X.\phi) &= \bigsqcup \{P \mid P \sqsubseteq \texttt{transformer}(P)\}
\end{align*}
where $a \Rightarrow b := \bigsqcup\{z \mid z \otimes a \sqsubseteq b\}$
is left residuation.
\end{definition}

\begin{theorem}[Environment Monotonicity with Polarity]
\label{thm:env-mono}
For formula $\phi$ and variable $i$ with polarity
$p \in \{\mathrm{true}, \mathrm{false}\}$:
if $\rho_1(j) = \rho_2(j)$ for all $j \neq i$ and
$\rho_1(i) \sqsubseteq \rho_2(i)$, then
\[
  p = \mathrm{true} \implies \texttt{qSatisfies}(\rho_1, \phi)
    \sqsubseteq \texttt{qSatisfies}(\rho_2, \phi)
\]
\[
  p = \mathrm{false} \implies \texttt{qSatisfies}(\rho_2, \phi)
    \sqsubseteq \texttt{qSatisfies}(\rho_1, \phi)
\]
\end{theorem}

\noindent
The proof (130~lines) proceeds by structural induction on~$\phi$,
handling each constructor: base cases ($\top, \bot, X$) are immediate;
negation flips polarity via antitonicity of residuation;
conjunction/disjunction preserve polarity via $\sqcap$/$\sqcup$
monotonicity; diamond/box use multiplication and residuation
monotonicity; and the fixed-point cases ($\mu$/$\nu$) shift the de
Bruijn variable index, extend the environment, and show every
pre/post-fixed point for $\rho_2$ is also one for $\rho_1$.

This enables the Knaster-Tarski theorem, proving that fixed points
exist and are computable as limits of approximation sequences
(\texttt{muApprox\_mono}, \texttt{nuApprox\_antimono}).

\begin{example}[Instantiations]
\begin{itemize}[leftmargin=1.5em]
  \item \textbf{Probabilistic}: $Q = [0,1]$ with standard multiplication.
  \item \textbf{PLN evidence}: $Q = \mathbb{R}_{\geq 0}^\infty \times \mathbb{R}_{\geq 0}^\infty$ (positive/negative evidence counts).
  \item \textbf{Fuzzy}: $Q = [0,1]^n$ (multi-valued truth).
  \item \textbf{Tropical}: $Q = \mathbb{R}_{\geq 0}^\infty$ with $\min$/$+$ operations.
\end{itemize}
\end{example}


%% ====================================================================
\section{Common Infrastructure}
\label{sec:common}
%% ====================================================================

The three process calculi share approximately 40\% of their
reduction and congruence infrastructure.  A shared library in
\texttt{ProcessCalculi/Common/} (4~files, 246~lines) eliminates this
duplication through typeclasses and generic lemmas.

\subsection{Typeclasses}

\texttt{Common/ProcessAlgebra.lean} (31~lines) defines:

\begin{lstlisting}[caption={Shared process algebra typeclasses}]
class HasPar (P : Type) where
  par : P -> P -> P

class HasNu (P : Type) where
  nu : P -> P

class HasNil (P : Type) where
  nil : P
\end{lstlisting}

\noindent
\texttt{Common/Congruence.lean} (64~lines) adds \texttt{HasSC} for
structural congruence and the generic pattern
\texttt{reduces\_modulo\_sc}.

\subsection{Reflexive-Transitive Closure}

\texttt{Common/Star.lean} (148~lines) provides two flavors of
reflexive-transitive closure:
\begin{itemize}[leftmargin=1.5em]
  \item \texttt{RTClosure}: Type-valued (carries witness data).
  \item \texttt{RTClosureProp}: Prop-valued (for logical reasoning).
\end{itemize}
Generic lifters---\texttt{liftUnary}, \texttt{liftBinLeft},
\texttt{liftBinRight}---allow multi-step reduction to be lifted
through any unary or binary context constructor.

\subsection{Instance Registration}

Each calculus registers instances:
\begin{itemize}[leftmargin=1.5em]
  \item \textbf{$\pi$-calculus}: \texttt{HasPar}, \texttt{HasNil},
    \texttt{HasSC} (wrapping the Type-valued SC in \texttt{Nonempty}).
  \item \textbf{$\rho$-calculus}: \texttt{HasSC} (Prop-valued SC
    directly).
  \item \textbf{MQ-calculus}: \texttt{HasPar}, \texttt{HasNil},
    \texttt{HasNu}, \texttt{HasSC} (all four typeclasses).
\end{itemize}

\noindent
All three share: structural congruence as an equivalence relation with
parallel commutativity/associativity and scope extrusion (unified via
\texttt{HasSC}); and the COMM rule shape (sender + receiver $\to$
result), differing only in payload consumption (substitution in
$\pi$/$\rho$, measurement in MQ).


%% ====================================================================
\section{Summary and Future Work}
\label{sec:conclusion}
%% ====================================================================

\subsection{Summary}

We have presented a comprehensive Lean~4.28 formalization of five
process-algebraic systems and three cross-calculus bridges, totaling
17{,}821 lines with zero \texttt{sorry} and zero custom axioms
(Table~\ref{tab:summary}).

Key results include:
\begin{enumerate}[leftmargin=1.5em]
  \item The Galois connection in $\rho$-calculus \emph{is} modal
    diamond-box duality (Theorem~\ref{thm:galois}), validating the OSLF
    construction.
  \item Weak operational correspondence for restriction-free
    $\pi$-processes (\texttt{calculus\_weak\_correspondence\_full\_encode}):
    forward simulation (Theorem~\ref{thm:forward-sim}), backward admin
    reflection with three-branch decomposition (reflexive,
    admin-canonical, or reflected $\pi$-step), and their composition
    into a single theorem---under restriction-freeness and Barendregt
    hygiene hypotheses, totaling 2{,}745~lines.
  \item Constructive both-outcomes witness for MQ measurement branching,
    with backend-parametric denotational semantics and post-collapse
    normalization.
  \item We conjecture that $\rho$-calculus is strictly more expressive
    than $\mu$-calculus (Conjecture~\ref{conj:rho-mu-expressiveness}):
    reflection cannot be captured by behavioral observation.
  \item Complete quantale-valued $\mu$-calculus semantics with
    environment monotonicity (Theorem~\ref{thm:env-mono}), enabling
    Knaster-Tarski fixed-point computation across probabilistic, fuzzy,
    tropical, and PLN-evidence domains.
  \item Spice calculus $n$-step lookahead is exactly iterated diamond
    $\langle\rangle^n$, proving that precognitive agents compute
    $\mu$-calculus temporal properties.
  \item Classical Hennessy-Milner equivalence (observational
    indistinguishability iff bisimilarity) proved assumption-free for
    finite-carrier SC-quotiented canonical $\rho$-relations, and under
    predecessor-finiteness for the global relation.  Forward
    image-finiteness is auto-discharged for all
    \texttt{LanguageDef}-based relations.  A structural sufficient condition
    for predecessor-finiteness (\texttt{LangPredFiniteSafe}~\cite{leanOSLF2025})
    identifies exactly which languages admit the global HM~iff without
    quotienting; the $\rho$-calculus is provably excluded.
  \item MeTTa-calculus with executable COMM/REFL reduction, premise-IR
    adequacy bridge, and a shared-fragment translation into
    $\rho$-calculus with forward/backward shared-core simulation and
    bisimulation at both step and step$^*$ levels
    (Section~\ref{sec:metta-calc}).
\end{enumerate}

\subsection{Future Work}

\begin{enumerate}[leftmargin=1.5em]
  \item \textbf{Unrestricted $\pi$-to-$\rho$ bisimulation}: extend the
    weak correspondence from restriction-free processes to full
    $\pi$-calculus (including $\nu$ and replication), which requires
    encoding name-server apparatus into the theorem statement.  Remove
    Barendregt side conditions via explicit alpha-renaming.  Package a
    single assumption-minimal weak bisimulation iff.
  \item \textbf{Global Hennessy-Milner for $\rho$-calculus}: lift the
    finite-carrier SC-quotiented HM~iff to the global unquotiented
    $\rho$-relation.  A decidable sufficient condition for
    predecessor-finiteness (\texttt{LangPredFiniteSafe}~\cite{leanOSLF2025})
    now provides the positive direction, but the $\rho$-calculus provably
    fails it (COMM erases the channel name).  The remaining path is
    quotienting to a finite-image semantics or proving an alternative
    finiteness argument that accounts for name erasure.
  \item \textbf{Full MeTTa$\leftrightarrow$$\rho$ bisimulation}: extend
    the shared-core (quote/drop/parallel) forward/backward simulation to
    \textsc{Comm} and \textsc{Refl} fragments; remove the
    source-anchored restriction by proving translation injectivity.
  \item \textbf{Multi-qubit MQ gates}: extend gate semantics beyond
    single-wire action; prove collapse preservation under composition
    and measurement idempotence.
  \item \textbf{PLN temporal bridge}: prove that PLN temporal operators
    correspond to modal formulas over a time-indexed LTS.
  \item \textbf{OSLF integration for MQ}: instantiate MQ processes as an
    OSLF language instance with $\Diamond \dashv \Box$ Galois
    connection.
  \item \textbf{Spice-MQ bridge}: connect $\rho$-calculus Spice
    $n$-step lookahead to MQ measurement prediction.
\end{enumerate}

\bibliographystyle{plain}
\bibliography{references}

\end{document}
